\documentclass[10pt,a4paper]{amsart}
\usepackage[margin=19mm]{geometry}
\usepackage{amsmath,amssymb,mathtools}
\usepackage[scr=rsfs,cal=euler]{mathalfa}
\usepackage{xfrac}
\usepackage[unicode,hidelinks,bookmarks=false]{hyperref}
\newcommand{\ie}{i.e.\ }
\newcommand{\cf}{cf.\ }
\newcommand{\resp}{respectively\ }
\newcommand{\Subsection}{\subsection}
\providecommand{\nobreakdash}{\nobreak\hbox{-}}
\setlength{\emergencystretch}{2em}
\multlinegap0pt
\usepackage{bm}
\usepackage{bbm}
\usepackage{dsfont}
\usepackage{xspace}
\usepackage{cancel}
\usepackage{mathtools}
\usepackage{amsthm}
\makeatletter
\@ifundefined{mytheo}{\newtheorem{mytheo}{Mytheo}}{}
\@ifundefined{myproposition}{\newtheorem{myproposition}[mytheo]{Proposition}}{}
\makeatother
\title[Completeness Relation in Renormalized Quantum Systems]{Completeness Relation in Renormalized Quantum Systems}
\author{Fatih Erman, O. Teoman Turgut}
\date{Source: arXiv:2409.05372v4 (CC BY 4.0)}
\begin{document}
\maketitle

\noindent\emph{Source and licence.} Completeness Relation in Renormalized Quantum Systems. Version arXiv:2409.05372v4 (CC BY 4.0), distributed under the Creative Commons Attribution 4.0 International licence, as recorded in the arXiv licence metadata for this exact version and shown on the abstract page https://arxiv.org/abs/2409.05372v4. Full rights notice: licenses/arxiv-2409.05372-CC-BY-4.0.txt.\par\smallskip

\noindent\textit{Editorial scope.} Counted unit: the Myproposition whose source label is \texttt{None} in the article (source0.tex lines 430--442, source label \texttt{None}), delivered together with the article's own complete original proof of it (source0.tex lines 444--481, one \texttt{proof} environment of 38 lines). No mathematics is written, repaired or re-derived: statement and proof are the article's own text line for line. Required context retained uncounted: greenfuncexpansion (lines 267--269); boundstatewavefunctionexact (lines 417--419); derivativePhi (lines 422--424). Every definition, display and label of those lines is retained unchanged. Omitted from this package and not redistributed: 1--266, 270--416, 420--421, 425--429, 443--443, 482--1015. Nothing inside a retained excerpt is omitted or altered: every retained body is delivered exactly as the article's lines are written, so no omission receipt of any kind accompanies this package. Citations: the retained excerpts carry no citation command at all, so no bibliography entry of the article is transcribed and no reference is redistributed. Reference adaptation: none was needed and none was made; every reference inside the delivered unit resolves inside it, so no unresolved reference or citation is produced. Printed slips kept literal: no printed word, symbol, hypothesis or number of the counted statement or of its proof was altered. Layout and class: the article's own class is \texttt{article} with packages including graphicx, epsfig, amssymb, amsmath, xcolor, amsthm, geometry, tikz, subcaption, caption, adjustbox, tocbibind, dcolumn, bm; the wrapper is the lane's pinned \texttt{amsart} editorial wrapper at 10pt on A4, which restates the theorem-like environments the retained text needs and adds the article's own macro definitions that the retained text needs (none), with extra wrapper packages (bm, bbm, dsfont, xspace, cancel, mathtools). The delivered numbering is the unit's own. Assets: no retained span contains a figure environment, a graphics-inclusion command, a PDF-page inclusion, a table or a TikZ picture; no image file is copied into this package, the package declares no asset dependency and no image file is redistributed with it. Licence: version arXiv:2409.05372v4 of this article is distributed under the Creative Commons Attribution 4.0 International licence, as recorded in the arXiv licence metadata for this exact version and shown on the abstract page https://arxiv.org/abs/2409.05372v4. A full rights notice is delivered at licenses/arxiv-2409.05372-CC-BY-4.0.txt. Characters: every retained excerpt is pure ASCII, so the delivered engine, XeLaTeX with its default Latin Modern text font, reports no missing character.
\begin{eqnarray}
    G_0(x,y|E) = \sum_{n=0}^{\infty} \frac{\phi_n(x) \overline{\phi_n(y)}}{E_n-E}  \;, \label{greenfuncexpansion}
\end{eqnarray}

\begin{eqnarray}
    \psi_k(x)=  \frac{G_0(x,a|E_{k}^{*})}{ \left(-\frac{d \Phi(E)}{d E}\bigg|_{E=E_{k}^{*}}\right)^{1/2}} \;. \label{boundstatewavefunctionexact}
\end{eqnarray}

\begin{equation}
    {d\Phi(E)\over dE}\Big|_{E_k^*}= - \sum_{n=0}^\infty { |\phi_n(a)|^2\over (E_n-E_k^*)^2} \;. \label{derivativePhi}
\end{equation}

\begin{myproposition} Let $\phi_n$ be orthonormal set of eigenfunctions of $H_0$, i.e., 

\begin{eqnarray}
H_0 \phi_n &=&E_n \phi_n \nonumber \\ \int_{\mathcal{M}} \overline{\phi_n(x)} \phi_m (x) \; d\mu(x) & =&  \delta_{nm}.
\end{eqnarray} 
%
Then, the eigenfunctions $\psi_n$ of H, which is formally $H_0$ modified by a delta interaction supported at $x=a$ are orthonormal, that is, 
%
\begin{eqnarray}
    \int_{\mathcal{M}} \overline{\psi_n(x)} \psi_m (x) \; d\mu( x) = \delta_{nm} \;,
\end{eqnarray}
%
where $D=1,2,3$. \end{myproposition}

\begin{proof} We first prove for $D=2,3$, where the renormalization is needed to define point delta interactions properly. 

Using bilinear expansion (\ref{greenfuncexpansion}) of the Green's function of $H_0$ and the eigenfunction (\ref{boundstatewavefunctionexact}), we obtain
%
\begin{eqnarray}
    & &  \int_{\mathcal{M}} \overline{\psi_n(x)} \psi_m (x) \; d\mu( x) =  \int_{\mathcal{M}} \frac{\overline{G_0(x,a|E_{n}^{*})}}{ \left(-\frac{d\Phi(E)}{d E}\bigg|_{E=E_{n}^{*}}\right)^{1/2}} \frac{G_0(x,a|E_{m}^{*})}{\left(-\frac{d \Phi(E)}{d E}\bigg|_{E=E_{m}^{*}}\right)^{1/2}}  \; d\mu(x) \nonumber \\  & & = \frac{1}{\left(-\frac{d \Phi(E)}{d E}\bigg|_{E=E_{n}^{*}}\right)^{1/2} \left(-\frac{d \Phi(E)}{d E}\bigg|_{E=E_{m}^{*}}\right)^{1/2}}\int_{\mathcal{M}} \sum_k \frac{\phi_k(a) \overline{\phi_k(x)}}{E_k-E_{n}^*}  \sum_l \frac{\phi_l(x) \overline{\phi_l(a)}}{E_l-E_{m}^*}  \; d\mu(x) \;.
\end{eqnarray}
%
Interchanging the order of summation and integration and using the fact that $\phi_k$'s are orthonormal functions, we have
%
\begin{eqnarray}
 \int_{\mathcal{M}} \overline{\psi_n(x)} \psi_m (x) \; d\mu( x) = \frac{1}{\left(-\frac{d \Phi(E)}{d E}\bigg|_{E=E_{n}^{*}}\right)^{1/2} \left(-\frac{d \Phi(E)}{d E}\bigg|_{E=E_{m}^{*}}\right)^{1/2}} \sum_k \frac{|\phi_k(a)|^2}{(E_k-E_{n}^*)(E_k-E_{m}^*)} \;. \label{orth2}
\end{eqnarray}
%
If $n=m$ in \ref{orth2}, then it is easy to show that the new eigenfunctions $\psi_n$'s are automatically normalized thanks to the identity (\ref{derivativePhi}):
%
\begin{eqnarray}
    \int_{\mathcal{M}}  |\psi_n(x)|^2 d\mu(x) = -\frac{1}{\frac{d \Phi(E)}{d E}\bigg|_{E=E_{n}^{*}}} \sum_{k=0}^{\infty} 
    \frac{|\phi_k(a)|^2}{(E_k-E_n^{*})^2} = 1 \;.
\end{eqnarray}
%
%
For the case $n \neq m$, we first formally decompose the expression in the summation with a cut-off $N$ as a sum of two partial fractions
%
\begin{eqnarray}
    \sum_{k=0}^{N} \frac{|\phi_k(a)|^2}{(E_k-E_{n}^*)(E_k-E_{m}^*)} =  \sum_{k=0}^{N} \frac{|\phi_k(a)|^2}{(E_{n}^{*}-E_{m}^*)} \left( \frac{1}{E_k-E_{n}^*}-\frac{1}{E_k-E_{m}^*}\right) \;.
\end{eqnarray}
%
As explained in the renormalization procedure, each term $\sum_{k=0}^{N} \frac{|\phi_k(a)|^2}{E_k-E_{n}^*}$ is divergent as $N \to \infty$. Motivated by this, we add and subtract $\frac{1}{\alpha_R} + \sum_{k=0}^{N} \frac{|\phi_k(a)|^2}{E_k + \mu^2}$ to the above expression and obtain in the limit $N \to \infty$
%
\begin{eqnarray} \int_{\mathcal M} \overline{\psi_n(x)} \psi_m (x) \; d\mu( x) = \frac{1}{\left(E_{n}^* - E_{m}^* \right)} \frac{\left(\Phi(E_{n}^{*}) - \Phi(E_{m}^*)\right)}{\left(-\frac{d\Phi(E)}{d E}\bigg|_{E=E_{n}^{*}}\right)^{1/2} \left(-\frac{d \Phi(E)}{d E}\bigg|_{E=E_{m}^{*}}\right)^{1/2}} \;.
\end{eqnarray}
%
Since the zeroes of the function $\Phi$ are the bound state of the modified system, that is, $\Phi(E_{n}^{*})=0$ and $\Phi(E_{m}^{*})=0$ for all $n, m$ (when $n\neq m$), this completes our proof of the orthogonality of eigenfunctions for the modified  Hamiltonian having discrete spectrum.   


The case for $D=1$ can easily be proved by following the same steps introduced above, except that there is no need for renormalization. 
\end{proof}

\end{document}
